\documentclass[12pt,a4paper]{article}
\usepackage{../../../myconfig}

\begin{document}

\titlebox{Chủ đề: Oxyz}{Phương trình mặt phẳng}{Oxyz: Phương trình mặt phẳng}

\questionLinesManual{0}{1}{%
Trong không gian với hệ trục tọa độ $Oxyz$, cho mặt phẳng $(P)$ có phương trình $x - 3y - z + 8 = 0$. Vectơ nào sau đây là một vectơ pháp tuyến của $(P)$?
}{%
\begin{tasks}(4)
    \task $\vec{n}_1 = (1; -3; 1)$
    \task $\vec{n}_2 = (1; -3; -1)$
    \task $\vec{n}_3 = (1; -3; 8)$
    \task $\vec{n}_4 = (1; 3; 8)$
\end{tasks}
}

% --- CÂU 2: SỬA ĐIỂM N ĐỂ CÓ ĐÁP ÁN B ĐÚNG ---
\questionLinesManual{4}{2}{%
Trong không gian $Oxyz$, cho mặt phẳng $(\alpha)\colon 2x - 3y - z + 2 = 0$. Điểm nào dưới đây nằm trên mặt phẳng $(\alpha)$?
}{%
\begin{tasks}(4)
    \task $K(1; 0; 1)$
    \task $N(3; 4; -4)$
    \task $M(-1; 3; 0)$
    \task $P(2; -3; 1)$
\end{tasks}
}

\questionLinesManual{4}{3}{%
Trong không gian $Oxyz$, phương trình nào sau đây là phương trình của mặt phẳng $(Oyz)$?
}{%
\begin{tasks}(4)
    \task $x + 1 = 0$
    \task $y = 0$
    \task $x = 0$
    \task $z = 0$
\end{tasks}
}

\questionLinesManual{6}{4}{%
Trong không gian $Oxyz$, phương trình mặt phẳng đi qua gốc tọa độ $O(0;0;0)$ và nhận véctơ $\vec{n} = (2; 0; 1)$ làm véctơ pháp tuyến là:
}{%
\begin{tasks}(4)
    \task $x + 2y = 0$
    \task $2x + z = 0$
    \task $x + 2z = 0$
    \task $2x + y = 0$
\end{tasks}
}

\questionLinesManual{6}{5}{%
Trong không gian $Oxyz$, phương trình mặt phẳng đi qua điểm $A(1; 2; -3)$ và nhận $\vec{n} = (1; -2; 3)$ làm véctơ pháp tuyến là:
}{%
\begin{tasks}(2)
    \task $x - 2y + 3z + 12 = 0$
    \task $x - 2y + 3z - 12 = 0$
    \task $x - 2y - 3z - 6 = 0$
    \task $x - 2y - 3z + 6 = 0$
\end{tasks}
}

\questionLinesManual{6}{6}{%
Trong không gian $Oxyz$, cho điểm $M(1; 2; -1)$ và mặt phẳng $(P)\colon x + 2y + z = 0$. Mặt phẳng $(Q)$ đi qua $M$ và song song với $(P)$ có phương trình là:
}{%
\begin{tasks}(2)
    \task $x + 2y + z + 4 = 0$
    \task $x + 2y + z - 1 = 0$
    \task $x + 2y + z - 4 = 0$
    \task $x + 2y - z - 6 = 0$
\end{tasks}
}

\questionLinesManual{6}{7}{%
Trong không gian $Oxyz$, cho hai điểm $A(5; -4; 2)$ và $B(1; 2; 4)$. Mặt phẳng đi qua $A$ và vuông góc với đoạn thẳng $AB$ có phương trình là:
}{%
\begin{tasks}(2)
    \task $2x - 3y - z - 20 = 0$
    \task $3x - y + 3z - 25 = 0$
    \task $2x - 3y - z + 8 = 0$
    \task $3x - y + 3z - 13 = 0$
\end{tasks}
}

\questionLinesManual{6}{8}{%
Trong không gian $Oxyz$, cho hai điểm $A(-1; 2; 0)$ và $B(3; 0; 2)$. Mặt phẳng trung trực của đoạn thẳng $AB$ có phương trình là:
}{%
\begin{tasks}(2)
    \task $x + y + z - 3 = 0$
    \task $2x - y + z + 2 = 0$
    \task $2x + y + z - 4 = 0$
    \task $2x - y + z - 2 = 0$
\end{tasks}
}

\questionLinesManual{6}{9}{%
Trong không gian $Oxyz$, mặt phẳng đi qua ba điểm $A(4; 0; 0)$, $B(0; -2; 0)$ và $C(0; 0; 2)$ có phương trình là:
}{%
\begin{tasks}(2)
    \task $\dfrac{x}{4} + \dfrac{y}{2} + \dfrac{z}{-2} = 1$
    \task $\dfrac{x}{4} + \dfrac{y}{-2} + \dfrac{z}{2} = 0$
    \task $\dfrac{x}{-4} + \dfrac{y}{2} + \dfrac{z}{-2} = 1$
    \task $\dfrac{x}{4} + \dfrac{y}{-2} + \dfrac{z}{2} = 1$
\end{tasks}
}

% --- CÂU 10: SỬA TỌA ĐỘ 3 ĐIỂM ĐỂ RA ĐÚNG ĐÁP ÁN C ---
\questionLinesManual{8}{10}{%
Trong không gian $Oxyz$, cho ba điểm $A(1; 1; 0)$, $B(2; 3; 3)$ và $C(0; 1; 3)$. Mặt phẳng đi qua ba điểm $A, B, C$ có phương trình là:
}{%
\begin{tasks}(2)
    \task $x - y + 2z - 5 = 0$
    \task $x + 2y - 3z + 4 = 0$
    \task $3x - 3y + z = 0$
    \task $x + y - 2z + 3 = 0$
\end{tasks}
}

\questionLinesManual{12}{11}{%
Trong không gian $Oxyz$, cho hai điểm $A(0; 1; 0)$, $B(2; 0; 1)$ và mặt phẳng $(P)\colon x - y - 1 = 0$. Phương trình mặt phẳng đi qua hai điểm $A, B$ đồng thời vuông góc với $(P)$ là:
}{%
\begin{tasks}(2)
    \task $x + y - 3z - 1 = 0$
    \task $2x + 2y - 5z - 2 = 0$
    \task $x - 2y - 6z + 2 = 0$
    \task $x + y - z - 1 = 0$
\end{tasks}
}

\questionLinesManual{12}{12}{%
Trong không gian $Oxyz$, phương trình mặt phẳng đi qua điểm $A(1; 1; 1)$ đồng thời vuông góc với cả hai mặt phẳng $(\alpha)\colon x + y - z - 2 = 0$ và $(\beta)\colon x - y + z - 1 = 0$ là:
}{%
\begin{tasks}(2)
    \task $y + z - 2 = 0$
    \task $x + y + z - 3 = 0$
    \task $x - 2y + z = 0$
    \task $x + z - 2 = 0$
\end{tasks}
}

\questionLinesManual{12}{13}{%
Trong không gian $Oxyz$, mặt phẳng chứa trục $Ox$ và vuông góc với mặt phẳng $(\alpha)\colon 2x - 3y + 4z - 5 = 0$ có phương trình là:
}{%
\begin{tasks}(2)
    \task $3x + 2y = 0$
    \task $4y - 3z = 0$
    \task $4y + 3z = 0$
    \task $3x + 4y - 1 = 0$
\end{tasks}
}

% --- CÂU 14: SỬA CHỮ "1" THÀNH "x" Ở ĐẦU PHƯƠNG TRÌNH ---
\questionLinesManual{12}{14}{%
Trong không gian $Oxyz$, khoảng cách từ điểm $M(-1; 2; -3)$ đến mặt phẳng $(P)\colon x - 2y + 2z - 1 = 0$ bằng:
}{%
\begin{tasks}(4)
    \task $6$
    \task $4$
    \task $\dfrac{17}{3}$
    \task $17$
\end{tasks}
}

\questionLinesManual{12}{15}{%
Trong không gian $Oxyz$, khoảng cách giữa hai mặt phẳng song song $(P)\colon x + 2y + 2z - 10 = 0$ và $(Q)\colon x + 2y + 2z - 3 = 0$ bằng:
}{%
\begin{tasks}(4)
    \task $\dfrac{2}{3}$
    \task $\dfrac{7}{3}$
    \task $\dfrac{8}{3}$
    \task $\dfrac{4}{3}$
\end{tasks}
}

\questionLinesManual{12}{16}{%
Trong không gian $Oxyz$, hai mặt phẳng $(\alpha)\colon 3x - 2y - z + 5 = 0$ và $(\beta)\colon x - y + 5z - 3 = 0$ có vị trí tương đối là:
}{%
\begin{tasks}(2)
    \task Song song nhau
    \task Vuông góc nhau
    \task Cắt nhau không vuông góc
    \task Trùng nhau
\end{tasks}
}

\questionLinesManual{12}{17}{%
Trong không gian $Oxyz$, cho hai mặt phẳng $(P)\colon 2x + my + 3z - 5 = 0$ và $(Q)\colon nx - 8y - 6z + 2 = 0$. Tìm $m, n$ để $(P)$ song song với $(Q)$.
}{%
\begin{tasks}(2)
    \task $m = n = -4$
    \task $m = 4; n = -4$
    \task $m = -4; n = 4$
    \task $m = n = 4$
\end{tasks}
}

\questionLinesManual{12}{18}{%
Trong không gian $Oxyz$, cho mặt phẳng $(P)\colon 2x + 2y - z - 1 = 0$. Mặt phẳng nào sau đây song song với $(P)$ và cách $(P)$ một khoảng bằng $3$?
}{%
\begin{tasks}(2)
    \task $(Q)\colon 2x + 2y - z + 10 = 0$
    \task $(Q)\colon 2x + 2y - z + 4 = 0$
    \task $(Q)\colon 2x + 2y - z + 8 = 0$
    \task $(Q)\colon 2x + 2y - z - 8 = 0$
\end{tasks}
}

\questionLinesManual{12}{19}{%
Trong không gian $Oxyz$, cho điểm $M(1; 2; 3)$. Viết phương trình mặt phẳng $(P)$ đi qua $M$ và cắt các trục $Ox, Oy, Oz$ lần lượt tại $A, B, C$ sao cho $M$ là trọng tâm tam giác $ABC$.
}{%
\begin{tasks}(2)
    \task $6x + 3y + 2z + 18 = 0$
    \task $6x + 3y + 2z + 6 = 0$
    \task $6x + 3y + 2z - 18 = 0$
    \task $6x + 3y + 2z - 6 = 0$
\end{tasks}
}

\questionLinesManual{12}{20}{%
Trong không gian $Oxyz$, cho điểm $M(1; 2; 5)$. Mặt phẳng $(P)$ đi qua điểm $M$ cắt các trục tọa độ $Ox, Oy, Oz$ lần lượt tại $A, B, C$ sao cho $M$ là trực tâm tam giác $ABC$. Phương trình của mặt phẳng $(P)$ là:
}{%
\begin{tasks}(2)
    \task $x + 2y + 5z - 30 = 0$
    \task $x + y + z - 8 = 0$
    \task $\dfrac{x}{5} + \dfrac{y}{2} + \dfrac{z}{1} = 1$
    \task $5x + 2y + z - 14 = 0$
\end{tasks}
}

\end{document}